\chapter{Análisis de requisitos}
\label{chap:requisitos}

\ph{Este capítulo describe \textbf{qué} debe hacer el sistema, sin entrar todavía en cómo
lo hace. Es la base contractual del TFG: si un requisito no aparece aquí,
no debería darse por hecho que el sistema lo cumple.}

\section{Descripción del sistema}
\label{sec:descripcion_sistema}

\ph{Visión general del sistema en una página: qué es, qué resuelve y para quién.}

\ph{Cubre:}
\begin{itemize}
  \item \ph{Tipo de sistema (aplicación web, CLI, librería, servicio\dots).}
  \item \ph{Qué entradas acepta y qué salidas produce.}
  \item \ph{Subsistemas o módulos principales en los que se divide.}
  \item \ph{Actores que interactúan con él (anónimo, autenticado, administrador\dots).}
\end{itemize}

\ph{Ejemplo de subsistemas:}

\begin{itemize}
  \item \textbf{\ph{Subsistema 1}}: \ph{responsabilidad principal.}
  \item \textbf{\ph{Subsistema 2}}: \ph{responsabilidad principal.}
  \item \textbf{\ph{Subsistema 3}}: \ph{responsabilidad principal.}
\end{itemize}

\ph{Ejemplo de actores:}

\begin{itemize}
  \item \textbf{\ph{Actor 1}}: \ph{rol y nivel de acceso.}
  \item \textbf{\ph{Actor 2}}: \ph{rol y nivel de acceso.}
\end{itemize}

\section{Requisitos de información}
\label{sec:ri}

\ph{Describe las entidades de información que el sistema gestiona, con sus
atributos y restricciones (claves, unicidad, tipos, opcionalidad). Suele haber
una tabla por entidad.}

\subsection*{RI-01: \ph{Nombre de la entidad}}

\ph{Breve descripción del propósito de la entidad.}

\begin{table}[H]
\centering
\small
\begin{tabular}{@{}P{3.5cm} P{3.cm} P{7.2cm}@{}}
\toprule
\textbf{Atributo} & \textbf{Tipo} & \textbf{Descripción} \\
\midrule
\ph{atributo\_1} & \ph{Tipo} & \ph{Descripción y restricciones.} \\
\ph{atributo\_2} & \ph{Tipo} & \ph{Descripción y restricciones.} \\
\ph{atributo\_3} & \ph{Tipo} & \ph{Descripción y restricciones.} \\
\bottomrule
\end{tabular}
\caption{Atributos del requisito de información RI-01.}
\label{tab:ri01}
\end{table}

\subsection*{RI-02: \ph{Nombre de la entidad}}

\ph{Breve descripción del propósito de la entidad.}

\begin{table}[H]
\centering
\small
\begin{tabular}{@{}P{3cm} P{2.8cm} P{7.2cm}@{}}
\toprule
\textbf{Atributo} & \textbf{Tipo} & \textbf{Descripción} \\
\midrule
\ph{atributo\_1} & \ph{Tipo} & \ph{Descripción y restricciones.} \\
\ph{atributo\_2} & \ph{Tipo} & \ph{Descripción y restricciones.} \\
\bottomrule
\end{tabular}
\caption{Atributos del requisito de información RI-02.}
\label{tab:ri02}
\end{table}

\section{Requisitos funcionales}
\label{sec:rf}

\ph{Lista las capacidades concretas que el sistema debe ofrecer, agrupadas por
bloques temáticos (gestión de usuarios, lógica de negocio principal, exploración\dots).
Cada requisito tiene un identificador único, un nombre, una descripción
y una prioridad (Alta/Media/Baja).}

\begin{longtable}{P{1.7cm} P{3.6cm} P{7cm} P{1.5cm}}
\toprule
\textbf{ID} & \textbf{Nombre} & \textbf{Descripción} & \textbf{Prior.} \\
\midrule
\endhead
RF-01 & \ph{Nombre del RF}  & \ph{Descripción del comportamiento esperado.} & \ph{Alta}  \\[0.4em]
RF-02 & \ph{Nombre del RF}  & \ph{Descripción del comportamiento esperado.} & \ph{Alta}  \\[0.4em]
RF-03 & \ph{Nombre del RF}  & \ph{Descripción del comportamiento esperado.} & \ph{Media} \\[0.4em]
\midrule
RF-04 & \ph{Nombre del RF}  & \ph{Descripción del comportamiento esperado.} & \ph{Alta}  \\[0.4em]
RF-05 & \ph{Nombre del RF}  & \ph{Descripción del comportamiento esperado.} & \ph{Alta}  \\[0.4em]
\midrule
RF-06 & \ph{Nombre del RF}  & \ph{Descripción del comportamiento esperado.} & \ph{Media} \\
\bottomrule
\caption{Requisitos funcionales del sistema.}
\label{tab:rf}
\end{longtable}

\section{Requisitos no funcionales}
\label{sec:rnf}

\ph{Atributos transversales del sistema: seguridad, rendimiento, mantenibilidad,
calidad de código, despliegue, estándares\dots Son tan vinculantes como los
funcionales y deben poder verificarse. Formula cada RNF con una MÉTRICA y un
UMBRAL comprobables, no con adjetivos vagos: en lugar de «el sistema debe ser
rápido», escribe «el 95\% de las peticiones se atiende en menos de 200~ms».
El Capítulo~\ref{chap:resultados} usará estos umbrales como criterios de
éxito medibles, así que un RNF que no se pueda medir aquí no se podrá
evaluar allí. El umbral de cada RNF cuantitativo debe coincidir con el del
criterio de éxito asociado de la Tabla~\ref{tab:criterios-exito}: un único
umbral, dos vistas.}

\begin{longtable}{P{1.7cm} P{3.6cm} P{7cm} P{1.5cm}}
\toprule
\textbf{ID} & \textbf{Nombre} & \textbf{Descripción} & \textbf{Prior.} \\
\midrule
\endhead
RNF-01 & \ph{Nombre del RNF} & \ph{Atributo + métrica + umbral, p.\,ej.\ «el 95\% de las peticiones se atiende en menos de 200~ms, medido con una prueba de carga».} & \ph{Alta} \\[0.4em]
RNF-02 & \ph{Nombre del RNF} & \ph{Atributo de seguridad/rendimiento/calidad con métrica, umbral y forma de medirlo.} & \ph{Alta} \\[0.4em]
RNF-03 & \ph{Nombre del RNF} & \ph{Atributo de seguridad/rendimiento/calidad con métrica, umbral y forma de medirlo.} & \ph{Media} \\[0.4em]
\midrule
RNF-04 & \ph{Nombre del RNF} & \ph{Atributo de seguridad/rendimiento/calidad con métrica, umbral y forma de medirlo.} & \ph{Alta} \\[0.4em]
RNF-05 & \ph{Nombre del RNF} & \ph{Atributo de seguridad/rendimiento/calidad con métrica, umbral y forma de medirlo.} & \ph{Media} \\
\bottomrule
\caption{Requisitos no funcionales del sistema.}
\label{tab:rnf}
\end{longtable}

\section{Reglas de negocio}
\label{sec:reglas_negocio}

\ph{Las reglas de negocio (RN) son afirmaciones sobre cómo se comporta el
dominio que el sistema debe respetar, independientemente de la tecnología o
de la interfaz. No son requisitos funcionales (qué hace el sistema) ni
requisitos no funcionales (atributos transversales): son restricciones
lógicas del problema. Son las que se rompen últimas si cambian: no porque
hayamos refactorizado, sino porque el negocio ha cambiado.}

\begin{longtable}{P{1.7cm} P{4cm} P{8.8cm}}
\toprule
\textbf{ID} & \textbf{Nombre} & \textbf{Descripción} \\
\midrule
\endhead
RN-01 & \ph{Nombre de la regla} & \ph{Enunciado preciso de la regla y, si
   procede, la condición que la activa o la fuente que la justifica
   (normativa, decisión de producto, restricción del dominio).} \\[0.4em]
RN-02 & \ph{Nombre de la regla} & \ph{Enunciado preciso.} \\[0.4em]
RN-03 & \ph{Nombre de la regla} & \ph{Enunciado preciso.} \\
\bottomrule
\caption{Reglas de negocio del dominio.}
\label{tab:rn}
\end{longtable}

\section{Restricciones del proyecto}
\label{sec:restricciones}

\ph{Limitaciones externas que acotan el espacio de soluciones del proyecto.
A diferencia de los RNF, no son atributos del sistema construido sino
condicionantes del proceso de construcción.}

\begin{longtable}{P{1.7cm} P{3.5cm} P{9.3cm}}
\toprule
\textbf{ID} & \textbf{Tipo} & \textbf{Descripción} \\
\midrule
\endhead
RES-01 & \ph{Tecnológica} &
   \ph{p.\,ej.\ uso obligatorio de un lenguaje, framework o plataforma
   concretos.} \\[0.4em]
RES-02 & \ph{Legal / regulatoria} &
   \ph{p.\,ej.\ cumplimiento de RGPD, normativa sectorial, accesibilidad WCAG.} \\[0.4em]
RES-03 & \ph{Organizativa} &
   \ph{p.\,ej.\ uso de la infraestructura del grupo de investigación,
   integración con sistemas internos existentes.} \\[0.4em]
RES-04 & \ph{Económica / temporal} &
   \ph{p.\,ej.\ presupuesto cero, plazo fijado por el calendario académico.} \\
\bottomrule
\caption{Restricciones del proyecto.}
\label{tab:res}
\end{longtable}

\section{Casos de uso}
\label{sec:casos_uso}

\subsection{Diagrama general de casos de uso}
\label{subsec:cu_diagrama}

\ph{Resumen narrativo del diagrama: qué actores hay, qué casos de uso conectan
con cada actor y qué relaciones \texttt{include} o \texttt{extend} existen
entre casos de uso.}

\ph{(Espacio reservado para el diagrama de casos de uso del sistema, no incluido
en esta plantilla.)}

\subsection{Especificación de casos de uso principales}
\label{subsec:cu_especificacion}

\ph{Especifica con detalle los casos de uso más relevantes. Cada uno usa una ficha
con los campos: actores, descripción, precondiciones, flujo principal,
flujos alternativos y postcondiciones. Repite la ficha tantas veces como casos
de uso quieras especificar.}

\begin{table}[H]
\centering
\small
\renewcommand{\arraystretch}{1.35}
\begin{tabular}{@{}>{\bfseries\RaggedRight\arraybackslash}p{4cm} P{10cm}@{}}
\toprule
\multicolumn{2}{l}{\textbf{CU-01 --- \ph{Nombre del caso de uso}}} \\
\midrule
Actores        & \ph{Actor/es implicado/s.} \\
Descripción    & \ph{Resumen del objetivo del caso de uso.} \\
Precondiciones & \ph{Estado previo necesario para iniciar el caso de uso.} \\
Flujo principal &
  \ph{1. Paso 1.\newline
       2. Paso 2.\newline
       3. Paso 3.\newline
       4. Paso 4.\newline
       5. Paso 5.} \\
Flujo alternativo &
  \ph{Xa. Variante o error y respuesta del sistema.} \\
Postcondiciones & \ph{Estado del sistema tras la ejecución correcta.} \\
\bottomrule
\end{tabular}
\caption{Especificación del caso de uso CU-01.}
\label{tab:uc01}
\end{table}

\begin{table}[H]
\centering
\small
\renewcommand{\arraystretch}{1.35}
\begin{tabular}{@{}>{\bfseries\RaggedRight\arraybackslash}p{4cm} P{10cm}@{}}
\toprule
\multicolumn{2}{l}{\textbf{CU-02 --- \ph{Nombre del caso de uso}}} \\
\midrule
Actores        & \ph{Actor/es implicado/s.} \\
Descripción    & \ph{Resumen del objetivo del caso de uso.} \\
Precondiciones & \ph{Estado previo necesario.} \\
Flujo principal &
  \ph{1. Paso 1.\newline
       2. Paso 2.\newline
       3. Paso 3.} \\
Flujo alternativo & \ph{Xa. Variante o error.} \\
Postcondiciones & \ph{Estado tras la ejecución.} \\
\bottomrule
\end{tabular}
\caption{Especificación del caso de uso CU-02.}
\label{tab:uc02}
\end{table}

\begin{table}[H]
\centering
\small
\renewcommand{\arraystretch}{1.35}
\begin{tabular}{@{}>{\bfseries\RaggedRight\arraybackslash}p{4cm} P{10cm}@{}}
\toprule
\multicolumn{2}{l}{\textbf{CU-03 --- \ph{Nombre del caso de uso}}} \\
\midrule
Actores        & \ph{Actor/es.} \\
Descripción    & \ph{Resumen.} \\
Precondiciones & \ph{Estado previo.} \\
Flujo principal &
  \ph{1. Paso 1.\newline
       2. Paso 2.\newline
       3. Paso 3.\newline
       4. Paso 4.} \\
Flujo alternativo & \ph{Xa. Variante.} \\
Postcondiciones & \ph{Estado tras la ejecución.} \\
\bottomrule
\end{tabular}
\caption{Especificación del caso de uso CU-03.}
\label{tab:uc03}
\end{table}

\section{Matriz de trazabilidad}
\label{sec:trazabilidad}

\ph{La matriz de trazabilidad cruza los requisitos funcionales con los
casos de uso que los implementan y, opcionalmente, con las pruebas que los
verifican. Demuestra que ningún requisito queda huérfano y que cada caso
de uso responde a uno o más requisitos identificados.}

\ph{Marca con \texttimes\ la celda en la que el caso de uso (o la prueba)
cubre el requisito de la fila. Una fila vacía es un requisito sin cobertura;
una columna vacía es un caso de uso sin requisito que lo justifique.
Ambos son señales de alerta.}

\subsection{Requisitos funcionales \texorpdfstring{$\leftrightarrow$}{<->} casos de uso}
\label{subsec:trazabilidad_rf_cu}

\ph{Ajusta el número de columnas al número real de casos de uso
especificados (la tabla de ejemplo muestra seis columnas, pero solo se han
incluido tres fichas de ejemplo).}

\begin{longtable}{P{1.7cm} *{6}{P{1.4cm}}}
\toprule
\textbf{} & \textbf{CU-01} & \textbf{CU-02} & \textbf{CU-03} & \textbf{CU-04} & \textbf{CU-05} & \textbf{CU-06} \\
\midrule
\endhead
RF-01 & \ph{$\times$} &               &               &               &               &               \\
RF-02 &               & \ph{$\times$} &               &               &               &               \\
RF-03 &               &               & \ph{$\times$} &               &               &               \\
RF-04 &               &               &               & \ph{$\times$} &               &               \\
RF-05 &               &               &               &               & \ph{$\times$} &               \\
RF-06 &               &               &               &               &               & \ph{$\times$} \\
\bottomrule
\caption{Matriz de trazabilidad entre requisitos funcionales y casos de uso.}
\label{tab:trazabilidad_rf_cu}
\end{longtable}

\subsection{Requisitos \texorpdfstring{$\leftrightarrow$}{<->} pruebas}
\label{subsec:trazabilidad_rf_pruebas}

\ph{Cruce de los requisitos con la suite de pruebas (Capítulo~\ref{chap:pruebas}).
Cada requisito debería tener al menos una prueba que lo cubra. Indica el
identificador o nombre del test correspondiente.}

\begin{longtable}{P{2cm} P{4cm} P{8.5cm}}
\toprule
\textbf{Requisito} & \textbf{Tipo de prueba} & \textbf{Test / archivo} \\
\midrule
\endhead
RF-01 & \ph{Unitaria / Integración / E2E} & \ph{\texttt{test\_xxx::test\_caso}} \\[0.3em]
RF-02 & \ph{Unitaria / Integración / E2E} & \ph{\texttt{test\_xxx::test\_caso}} \\[0.3em]
RF-03 & \ph{Unitaria / Integración / E2E} & \ph{\texttt{test\_xxx::test\_caso}} \\[0.3em]
\midrule
RNF-01 & \ph{Tipo} & \ph{Test / verificación.} \\[0.3em]
RNF-02 & \ph{Tipo} & \ph{Test / verificación.} \\
\bottomrule
\caption{Matriz de trazabilidad entre requisitos y pruebas.}
\label{tab:trazabilidad_rf_pruebas}
\end{longtable}

% --- Transición al siguiente capítulo ---
\ph{Cierra el capítulo con un párrafo puente de 2--4 líneas: qué queda
establecido aquí y qué pregunta abre el capítulo siguiente. Por ejemplo:
«Con los requisitos especificados y trazados, queda fijado \emph{qué} debe
hacer el sistema. El Capítulo~\ref{chap:diseno} aborda \emph{cómo}
satisfacerlos: la arquitectura, los modelos y las decisiones de diseño que
dan forma a la solución».}
